\clearpage
\section*{Problem 1: Orthogonal Matching Pursuit (OMP)}
\addcontentsline{toc}{section}{Problem 1: Orthogonal Matching Pursuit}
\markboth{Problem 1}{Problem 1}

\textbf{Problem Statement:} 

\subsection*{Part 0: Preliminary Information and Data}

In part 2, use the matrix
$$A = \begin{bmatrix} 
1 & -1/2 & -1/2 \\
0 & \sqrt{3}/2 & -\sqrt{3}/2 \\
1 & 3 & 3
\end{bmatrix}$$

In part 3, use the matrix $B$ that is the $2 \times 3$ matrix consisting of the first two rows of $A$.

\subsection*{Singular Value Decomposition (SVD) and Pseudo-Inverse}

Given a matrix $M_{m \times n}$ with $\text{rank}(M) = r$, write
$$M = U\Sigma V^T$$
where $U_{m \times m}$ and $V_{n \times n}$ are orthogonal matrices, and $\Sigma_{m \times n}$ is a diagonal matrix with nonzero diagonal entries $\sigma_1 \geq \sigma_2 \geq \ldots \geq \sigma_r > 0$ called singular values of $M$.

\textbf{Definition.} Given $M$ as above, the pseudo-inverse $M^+$ is defined by
$$M^+ = V\Sigma^+ U^T$$
where $\Sigma^+_{n \times m}$ is the diagonal matrix whose non-zero entries are $\frac{1}{\sigma_1}, \frac{1}{\sigma_2}, \ldots, \frac{1}{\sigma_r}$.

\subsection*{Thresholding Operator}

The \textbf{hard thresholding operator} $H_s$ is defined as follows:
$$\mathbf{x} \to H_s(\mathbf{x}) \in \mathbb{R}^n$$

$H_s(\mathbf{x})$ sets all entries of $\mathbf{x}$ to zero except for $s$ largest (in absolute value) entries.

\subsection*{The Matrix $A_\mathcal{S}$}

Given a matrix $A$ and an index set $\mathcal{S}$, denote by $A_\mathcal{S}$ the matrix consisting of columns of $A$ with index in $\mathcal{S}$.

\noindent\rule{\textwidth}{0.4pt}

\vspace{1em}

\subsection*{Assignments and Questions}

\begin{enumerate}
\item Implement the pseudocode for the Orthogonal Matching Pursuit (OMP). The pseudocode is contained in a separate handwritten note attached to the Final project assignment. You may use the standard Matlab (or other language) commands to compute SVD and pseudo-inverse.

\item Run the code for $\mathbf{b} = A\mathbf{e}_1$ with $\mathbf{e}_1 = (1, 0, 0)^T$. See how well the obtained vector $\mathbf{x}$ matches $\mathbf{e}_1$.

\item Run the code again, this time using the matrix $B$ instead of $A$.

\item Compare the results obtained in parts 2, 3. Explain the connection with compressed sensing and sparse recovery.
\end{enumerate}

\vspace{1em}

\noindent\rule{\textwidth}{0.4pt}

\clearpage
\subsection*{Part 1: OMP Implementation}
\addcontentsline{toc}{subsection}{Part 1: OMP Implementation}

The Orthogonal Matching Pursuit algorithm was implemented following the provided pseudocode. Key components include:

\begin{itemize}
    \item Hard thresholding operator $H_s(\mathbf{x})$ that retains only the $s$ largest entries
    \item OMP iteration that computes correlations $\mathbf{g}_k = A^T \mathbf{r}_k$, applies hard thresholding, updates the support set, solves least squares using pseudo-inverse $A_\mathcal{S}^+$, and updates the residual
    \item Convergence detection based on residual norm threshold
\end{itemize}

See Appendix B for implementation details.

\vspace{1em}

\textbf{Matrix Properties:}

\textbf{Matrix A} ($3 \times 3$):
\begin{itemize}
    \item Rank: 3 (full rank)
    \item Condition number: 3.8842
    \item Singular values: $\sigma_1 = 4.3847$, $\sigma_2 = 1.2247$, $\sigma_3 = 1.1289$
    \item Square matrix providing $m=3$ measurements for $n=3$ unknowns
\end{itemize}

\textbf{Matrix B} ($2 \times 3$):
\begin{itemize}
    \item Rank: 2 (full row rank)
    \item Condition number: 1.0000
    \item Singular values: $\sigma_1 = 1.2247$, $\sigma_2 = 1.2247$
    \item Underdetermined system with $m=2$ measurements for $n=3$ unknowns
\end{itemize}

\vspace{2em}
\noindent\rule{\textwidth}{0.4pt}

\clearpage
\subsection*{Part 2: OMP with Matrix A}
\addcontentsline{toc}{subsection}{Part 2: OMP with Matrix A}

\textbf{Setup:} Target $\mathbf{e}_1 = (1, 0, 0)^T$, observation $\mathbf{b} = A\mathbf{e}_1 = (1, 0, 1)^T$

\textbf{Results for varying sparsity levels:}

\begin{table}[h]
\centering
\begin{tabular}{cccccc}
\toprule
$s$ & Recovered $\mathbf{x}$ & $\|\mathbf{x} - \mathbf{e}_1\|_2$ & $\|A\mathbf{x} - \mathbf{b}\|_2$ & Recovery & Feasible \\
\midrule
1 & $(1, 0, 0)^T$ & $3.33 \times 10^{-16}$ & $4.71 \times 10^{-16}$ & $\checkmark$ & $\checkmark$ \\
2 & $(1, 0, 0)^T$ & $3.33 \times 10^{-16}$ & $4.71 \times 10^{-16}$ & $\checkmark$ & $\checkmark$ \\
3 & $(1, 0, 0)^T$ & $3.33 \times 10^{-16}$ & $4.71 \times 10^{-16}$ & $\checkmark$ & $\checkmark$ \\
\bottomrule
\end{tabular}
\caption{OMP results for Matrix A with different sparsity levels}
\end{table}

\textbf{Analysis:} OMP achieves perfect recovery at $s=1$ because $\mathbf{e}_1$ is 1-sparse. The algorithm correctly identifies that column 1 has the highest normalized correlation with the observation vector. Both feasibility ($A\mathbf{x} = \mathbf{b}$) and recovery ($\mathbf{x} = \mathbf{e}_1$) are satisfied for all $s \geq 1$.

\vspace{2em}
\noindent\rule{\textwidth}{0.4pt}

\clearpage
\subsection*{Part 3: OMP with Matrix B}
\addcontentsline{toc}{subsection}{Part 3: OMP with Matrix B}

\textbf{Setup:} Target $\mathbf{e}_1 = (1, 0, 0)^T$, observation $\mathbf{b} = B\mathbf{e}_1 = (1, 0)^T$

\textbf{Results for varying sparsity levels:}

\begin{table}[h]
\centering
\begin{tabular}{cccccc}
\toprule
$s$ & Recovered $\mathbf{x}$ & $\|\mathbf{x} - \mathbf{e}_1\|_2$ & $\|B\mathbf{x} - \mathbf{b}\|_2$ & Recovery & Feasible \\
\midrule
1 & $(1, 0, 0)^T$ & $0.00 \times 10^{0}$ & $0.00 \times 10^{0}$ & $\checkmark$ & $\checkmark$ \\
2 & $(1, 0, 0)^T$ & $0.00 \times 10^{0}$ & $0.00 \times 10^{0}$ & $\checkmark$ & $\checkmark$ \\
\bottomrule
\end{tabular}
\caption{OMP results for Matrix B with different sparsity levels}
\end{table}

\textbf{Analysis:} OMP achieves perfect recovery at $s=1$ because $\mathbf{b}_B = (1, 0)^T$ has sparsity 1 in the column space—it requires only 1 column of $B$ for exact representation.

\vspace{2em}
\noindent\rule{\textwidth}{0.4pt}

\clearpage
\subsection*{Part 4: Comparison and Connection to Compressed Sensing}
\addcontentsline{toc}{subsection}{Part 4: Comparison}

\begin{table}[h]
\centering
\begin{tabular}{lcc}
\toprule
Property & Matrix A & Matrix B \\
\midrule
System type & Determined $(m=n)$ & Underdetermined $(m<n)$ \\
Dimensions & $3 \times 3$ & $2 \times 3$ \\
Rank & 3 (full) & 2 (full row rank) \\
Condition number & 3.8842 & 1.0000 \\
Sparsity of $\mathbf{e}_1$ & 1 & 1 \\
Min. $s$ for recovery & 1 & 1 \\
Recovery success & $\checkmark$ & $\checkmark$ \\
\bottomrule
\end{tabular}
\caption{Comparison of OMP performance on matrices A and B}
\end{table}

\subsubsection*{Key Observations}

\textbf{Feasibility vs. Recovery:}
\begin{itemize}
    \item \textbf{Matrix A:} Perfect recovery at $s=1$ for the 1-sparse signal $\mathbf{e}_1$
    \item \textbf{Matrix B:} Perfect recovery at $s=1$ for the 1-sparse signal $\mathbf{e}_1$
    \item Both matrices achieve feasibility (satisfy $\mathbf{Ax} = \mathbf{b}$) and recovery (recover $\mathbf{x} = \mathbf{e}_1$) when $s \geq$ sparsity$(\mathbf{e}_1)$
\end{itemize}

\textbf{Critical Insight: Normalized Correlations}

The success of OMP critically depends on using \emph{normalized correlations} when selecting columns:
$$\text{correlation}_j = \frac{|\mathbf{a}_j^T \mathbf{r}|}{\|\mathbf{a}_j\|_2 \cdot \|\mathbf{r}\|_2}$$

Without normalization, the algorithm selects based on dot product magnitude $|\mathbf{a}_j^T \mathbf{r}|$, which can be dominated by columns with large norms rather than those best aligned with the residual. This normalization ensures that OMP selects the column with the highest cosine similarity to the residual.

\textbf{Compressed Sensing Demonstration:}

Matrix B ($2 \times 3$, underdetermined) perfectly recovers the 1-sparse signal $\mathbf{e}_1$ from only $m=2$ measurements for $n=3$ unknowns. This demonstrates the fundamental principle of compressed sensing:

\begin{itemize}
    \item Sparse signals can be recovered from $m < n$ measurements
    \item Recovery succeeds when $s \geq$ sparsity of the signal
    \item The measurement matrix must satisfy certain properties (e.g., low coherence, RIP)
    \item Proper algorithm implementation (normalized correlations) is essential
\end{itemize}

\textbf{Conclusion:} Both matrices successfully recover the 1-sparse signal at $s=1$, validating the OMP algorithm. Matrix B's success with $m < n$ confirms that compressed sensing enables recovery from underdetermined systems when signals are sufficiently sparse. The key requirements are: (1) sparsity budget $s \geq$ true sparsity, (2) well-conditioned measurement matrix, and (3) normalized correlation-based column selection.

\vspace{2em}
\noindent\rule{\textwidth}{0.4pt}
